\documentclass[11pt]{article}
\usepackage[margin=1in]{geometry}
\usepackage{enumitem}
% =============================================================================
% Preambles/header.tex — shared LaTeX/Beamer preamble for this project
%
% Use from a Beamer lecture by prepending:
%     \input{header}
% and compiling with TEXINPUTS=../Preambles:$TEXINPUTS (the /compile-latex
% skill does this automatically).
%
% PALETTE CONTRACT. Color names defined here MUST match the names in
% Quarto/theme-template.scss so Beamer and Quarto renderings use the same
% palette. The scripts/check-palette-sync.sh script enforces this.
%
% To customize: edit the HEX values below AND the matching $variable in
% Quarto/theme-template.scss, then run scripts/check-palette-sync.sh.
% =============================================================================

% -----------------------------------------------------------------------------
% Engine detection + encoding
%
% Our /compile-latex skill uses XeLaTeX, where inputenc is unnecessary and
% can produce encoding warnings. Only load inputenc under pdfTeX.
% -----------------------------------------------------------------------------
\usepackage{iftex}
\ifPDFTeX
  \usepackage[utf8]{inputenc}
\fi

\usepackage{amsmath, amssymb, amsthm}
\usepackage{booktabs}
\usepackage{graphicx}

% -----------------------------------------------------------------------------
% PALETTE — mirrors Quarto/theme-template.scss variable names.
% Load xcolor and define colors BEFORE anything that references them
% (notably hyperref, hypersetup, and the Beamer color assignments).
% -----------------------------------------------------------------------------
\usepackage{xcolor}

\definecolor{primary-blue}{HTML}{012169}      % headings, accents
\definecolor{primary-gold}{HTML}{B9975B}      % emphasis, borders
\definecolor{highlight-yellow}{HTML}{F2A900}  % markers, alerts
\definecolor{light-bg}{HTML}{E8EDF5}          % title / section backgrounds
\definecolor{jet}{HTML}{1A1A1A}               % body text

% Semantic colors (used in diagrams and annotations)
\definecolor{positive}{HTML}{15803D}          % good / observed / identified
\definecolor{negative}{HTML}{B91C1C}          % bad / problematic / bias
\definecolor{neutral}{HTML}{525252}           % reference / context

% Highlight accents (match .hi-* classes in the SCSS)
\definecolor{hi-slate}{HTML}{314F4F}
\definecolor{hi-green}{HTML}{2E8B57}
\definecolor{hi-red}{HTML}{C41E3A}

% Semantic aliases so skill templates and snippets can write
% \textcolor{accent}{...} without hard-coding "primary-blue".
\colorlet{accent}{primary-blue}
\colorlet{accent2}{primary-gold}

% -----------------------------------------------------------------------------
% Hyperref — loaded AFTER xcolor + palette so link colors resolve.
% -----------------------------------------------------------------------------
\usepackage{hyperref}
\hypersetup{colorlinks=true, linkcolor=primary-blue, urlcolor=primary-blue,
            citecolor=primary-blue, pdfborder={0 0 0}}

% -----------------------------------------------------------------------------
% Course code — set once per deck (after \input{header}, before \begin{document})
% via \coursecode{CS401}. Shown in the footer on every slide so a deck's course
% is identifiable without opening the title page. Empty by default.
% -----------------------------------------------------------------------------
\makeatletter
\newcommand{\coursecode}[1]{\gdef\@coursecodetext{#1}}
\gdef\@coursecodetext{}
\makeatother

% -----------------------------------------------------------------------------
% Beamer theme assignments (only apply under Beamer).
%
% We detect Beamer via \@ifclassloaded{beamer}, which is the documented,
% non-internal way. This must be wrapped in \makeatletter/\makeatother
% because \@ifclassloaded contains an @-catcode character.
% -----------------------------------------------------------------------------
\makeatletter
\@ifclassloaded{beamer}{%
  \usefonttheme{professionalfonts}
  \setbeamercolor{structure}{fg=primary-blue}
  \setbeamercolor{frametitle}{fg=primary-blue}
  \setbeamercolor{title}{fg=primary-blue}
  \setbeamercolor{subtitle}{fg=primary-gold}
  \setbeamercolor{author}{fg=primary-blue}
  \setbeamercolor{institute}{fg=neutral}
  \setbeamercolor{date}{fg=primary-gold}
  \setbeamercolor{itemize item}{fg=primary-blue}
  \setbeamercolor{itemize subitem}{fg=primary-gold}
  \setbeamercolor{alerted text}{fg=primary-gold}
  \setbeamercolor{block title}{fg=white, bg=primary-blue}
  \setbeamercolor{block body}{bg=light-bg}
  % Semantic block variants: block=definition/concept (blue), exampleblock=
  % worked example (green/positive), alertblock=key takeaway or caution (gold).
  % Keep to <=2 colored boxes per slide across ALL three variants (INV-7).
  \setbeamercolor{block title example}{fg=white, bg=positive}
  \setbeamercolor{block body example}{bg=light-bg}
  \setbeamercolor{block title alerted}{fg=white, bg=primary-gold}
  \setbeamercolor{block body alerted}{bg=light-bg}

  % Minimal footer: course code (left) + page X / N (right), no navigation clutter
  \setbeamertemplate{navigation symbols}{}
  \setbeamertemplate{footline}{%
    \color{neutral}\footnotesize\hspace{0.7em}\@coursecodetext\hfill
    \insertframenumber/\inserttotalframenumber\hspace{0.7em}\vspace{0.3em}}
}{}
\makeatother

% -----------------------------------------------------------------------------
% TikZ setup — libraries the snippet gallery uses
% -----------------------------------------------------------------------------
\usepackage{tikz}
\usetikzlibrary{arrows.meta, positioning, calc, decorations.pathreplacing,
                fit, shapes.geometric, backgrounds}

% Shared TikZ styles — snippets and hand-written diagrams can pick these up
% via `\begin{tikzpicture}[every node/.style={...}]` or by naming specific
% styles (e.g., `\node[dag-node] (X) at (0,0) {X};`).
\tikzset{
  % Node styles for causal DAGs and similar
  dag-node/.style = {
    draw=primary-blue, fill=light-bg, rounded corners=2pt,
    minimum width=1.2cm, minimum height=0.9cm, align=center, font=\small
  },
  decision-node/.style = {
    draw=primary-gold, fill=white, diamond, aspect=1.8,
    minimum width=1.8cm, minimum height=1.2cm, align=center, font=\small,
    inner sep=1pt
  },
  % Edge styles — observed vs counterfactual
  observed-edge/.style = {-{Latex[length=2.5mm]}, thick, draw=jet},
  counterfactual-edge/.style = {-{Latex[length=2.5mm]}, thick, draw=neutral, dashed},
  confound-edge/.style = {-{Latex[length=2.5mm]}, thick, draw=negative, dashed},
  % Data-point styles for plots
  observed-dot/.style = {fill=primary-blue, draw=primary-blue, circle, inner sep=1.2pt},
  counterfactual-dot/.style = {fill=white, draw=neutral, circle, inner sep=1.2pt}
}

% -----------------------------------------------------------------------------
% Convenience macros
% -----------------------------------------------------------------------------
\newcommand{\muted}[1]{\textcolor{neutral}{#1}}
\newcommand{\key}[1]{\textcolor{primary-gold}{\textbf{#1}}}
\newcommand{\good}[1]{\textcolor{positive}{#1}}
\newcommand{\bad}[1]{\textcolor{negative}{#1}}

% Standout frame macro: full-bleed dark background for section transitions.
% Beamer-only (uses \begin{frame}/\setbeamercolor/\usebeamerfont) -- do not
% call from an article-class file (e.g. Notes/<CODE>/*-notes.tex). \muted,
% \key, \good, \bad, and \coursecode above are plain xcolor/gdef and are
% safe to use from either class; verified when Notes/article-class support
% was added.
% Expand: \transitionslide{Your transition title}
\newcommand{\transitionslide}[1]{%
  \begingroup
  \setbeamercolor{background canvas}{bg=primary-blue}
  \begin{frame}[plain]
    \centering
    \vfill
    {\usebeamerfont{title}\color{white}\Large #1\par}
    \vfill
  \end{frame}
  \endgroup
}

% Section divider frame (Beamer-only), an alternative to \transitionslide with a
% darker two-line style: near-black (jet) background, a small white label line
% above a larger gold title, and a thin gold rule along the bottom edge.
% Expand: \sectiondivider{Section 2}{Pointers}
\newcommand{\sectiondivider}[2]{%
  \begingroup
  \setbeamercolor{background canvas}{bg=jet}
  \begin{frame}[plain]
    \vspace*{3.0cm}
    {\color{white}\ttfamily\large #1\par}
    \vspace{0.5cm}
    {\color{primary-gold}\LARGE #2\par}
    \begin{tikzpicture}[remember picture, overlay]
      \draw[primary-gold, line width=2.5pt]
        ($(current page.south west)+(0,0.1)$) -- ($(current page.south east)+(0,0.1)$);
    \end{tikzpicture}
  \end{frame}
  \endgroup
}

% -----------------------------------------------------------------------------
% Channel watermark — "Concepts That Click" (YouTube).
%
% Article-class files (CompetitiveExam Q/A sets, Notes, Assignments) get a
% light diagonal draft watermark on every page plus a centered footer naming
% the channel. Beamer decks get a subtle TikZ background watermark instead
% (the footline already carries the course code). Centralized here so every
% MCQ PDF is branded without per-file edits.
% -----------------------------------------------------------------------------
\makeatletter
\@ifclassloaded{article}{%
  \usepackage{draftwatermark}
  \SetWatermarkText{Concepts That Click}
  \SetWatermarkScale{1}
  \SetWatermarkFontSize{2cm}
  \SetWatermarkLightness{0.86}
  \SetWatermarkAngle{45}
  \usepackage{fancyhdr}
  \pagestyle{fancy}
  \fancyhf{}
  \fancyfoot[C]{\footnotesize\textcolor{neutral}{Concepts That Click~|~YouTube: \href{https://www.youtube.com/@concepts.thatclick}{@concepts.thatclick}}}
  \renewcommand{\headrulewidth}{0pt}
  \renewcommand{\footrulewidth}{0pt}
  \fancypagestyle{plain}{%
    \fancyhf{}%
    \fancyfoot[C]{\footnotesize\textcolor{neutral}{Concepts That Click~|~YouTube: \href{https://www.youtube.com/@concepts.thatclick}{@concepts.thatclick}}}%
    \renewcommand{\headrulewidth}{0pt}%
    \renewcommand{\footrulewidth}{0pt}%
  }
}{}
\@ifclassloaded{beamer}{%
  \setbeamertemplate{background}{%
    \begin{tikzpicture}[remember picture, overlay]
      \node[rotate=45, opacity=0.055, align=center]
        at (current page.center)
        {\Huge\textcolor{neutral}{Concepts That Click}};
    \end{tikzpicture}%
  }
}{}
\makeatother 
\coursecode{JKSSB}
\renewcommand{\thesection}{06.\arabic{section}}
\newtheorem{example}{Example}[section]
\newenvironment{definitionbox}[1]{%
  \par\noindent\textbf{Definition (#1).}\ \itshape
}{\par\vspace{0.3em}}
\title{CPU vs GPU \& Microprocessor --- Detailed Lecture Notes}
\author{JKSSB Computer Awareness}
\date{\today}

\begin{document}
\maketitle

\noindent\textbf{Video lesson:}
\href{https://youtu.be/Vaod5Ox7lzM}{CPU vs GPU: What's the Difference? Microprocessor \& Microcontroller Explained}

\section{Serial and parallel processing}
\textbf{Serial processing} means that one task is carried out at a time, in
sequence: the first operation finishes before the next begins. \textbf{Parallel
processing} means that several tasks or operations proceed at the same time.
A serial workload rewards a processor with a few very fast units; a parallel
workload rewards a processor with many units that can work simultaneously.

\begin{definitionbox}{Serial processing}
Carrying out one task at a time, in sequence, so that each step completes
before the next starts.
\end{definitionbox}

\begin{definitionbox}{Parallel processing}
Carrying out several tasks or operations at the same time.
\end{definitionbox}

The difference matters because it explains why two processors can both be
described as ``fast'' while being designed for opposite jobs. A processor that
excels at finishing a single chain of instructions quickly is a
latency-optimised design. A processor that excels at completing a very large
number of small independent operations per second is a
throughput-optimised design.

\section{CPU versus GPU}
The CPU (Central Processing Unit) is the general-purpose processor of a
computer. It has \textbf{few, powerful cores}, and each core is capable of
running a sequential instruction stream quickly. Its design goal is
\textbf{low latency}: to complete one task with the least possible delay. In
the office-desktop running example (EX-01), the CPU executes the operating
system, \mbox{MS Word}, and spreadsheet logic step by step.

The GPU (Graphics Processing Unit) is a processor built around
\textbf{many simpler cores}. Each core does a small amount of work, and the
design goal is \textbf{high throughput}: a very large number of independent
operations per second. This makes the GPU a natural fit for pixels and
graphics, and also for scientific and machine-learning workloads that can be
expressed as many parallel calculations.

\begin{center}
\begin{tabular}{p{0.20\textwidth}p{0.34\textwidth}p{0.34\textwidth}}
\toprule
\textbf{Feature} & \textbf{CPU} & \textbf{GPU}\\
\midrule
Cores & Few, powerful & Many, simple\\
Design goal & Low latency & High throughput\\
Best suited to & Sequential, general-purpose work & Parallel work (pixels, ML)\\
Core type & General-purpose & Specialised, simpler\\
Example use & EX-01 desktop logic & Image and graphics operations\\
\bottomrule
\end{tabular}
\end{center}

\begin{example}
An item asks which processor is better suited to an image-processing task
that repeats the same operation on millions of pixels. Because those pixel
operations are independent and can run in parallel, the GPU --- with its many
simple cores --- is the better fit. The CPU's few powerful cores would each
handle the chain of pixels one after another.
\end{example}

The most common near-neighbour error is to treat GPU cores as if they were
interchangeable general-purpose CPU cores. They are not. GPU cores are
simpler and optimised for parallel throughput, so they are not a full
replacement for the CPU's general-purpose cores.

\section{Cores, threads, and hyper-threading}
A \textbf{core} is a physical processing unit. A quad-core CPU contains four
such units. A \textbf{thread} is a stream of instructions that a core
executes. \textbf{Hyper-threading} lets a single physical core present two
logical threads, so the core can stay busy while one thread is waiting.

Two logical threads on one core are not two physical cores. A quad-core
processor with hyper-threading may report eight logical threads, but it still
has four physical cores. Clock speed and core count together, not thread
count alone, determine raw capacity.

\begin{definitionbox}{Hyper-threading}
A technique in which one physical core presents two logical threads to the
operating system, improving utilisation while a thread waits.
\end{definitionbox}

\section{Cache depth: L1, L2, and L3}
Cache is small, fast memory held close to the core. It is organised in levels.
\textbf{L1} is the smallest and fastest level and sits nearest the core.
\textbf{L2} is larger and a little slower than L1. \textbf{L3} is the largest
and slowest of the three and is typically shared across cores. The order of
speed and closeness is L1, then L2, then L3, then main memory.

\begin{center}
\begin{tabular}{p{0.12\textwidth}p{0.30\textwidth}p{0.42\textwidth}}
\toprule
\textbf{Level} & \textbf{Relative size} & \textbf{Relative speed and location}\\
\midrule
L1 & Smallest & Fastest; closest to the core\\
L2 & Larger & Slower than L1\\
L3 & Largest & Slowest of the three; often shared across cores\\
\bottomrule
\end{tabular}
\end{center}

\section{Clock speed and multi-core}
\textbf{Clock speed} is measured in \textbf{GHz} (gigahertz), which counts
billions of cycles per second. A higher clock speed generally means more work
completed per second, all else being equal. A \textbf{multi-core} processor
places several cores on one chip, which enables parallel work. However, clock
speed alone never decides real performance: the number of cores, the cache
size, the amount of memory, and the input/output speed all contribute.

\section{Microprocessor versus microcontroller}
A \textbf{microprocessor} is a central processing unit implemented on a single
chip. It is a general-purpose processor, and software decides what it does.
On its own it still needs external memory and input/output devices to form a
complete working system. The EX-01 office desktop is built around a
microprocessor.

A \textbf{microcontroller} integrates more on the same chip: a microprocessor
\textbf{plus} memory \textbf{plus} input/output peripherals. It therefore
forms a near-complete computer on one chip and is used for
\textbf{embedded/special-purpose} control, for example in washing machines,
remote controls, and appliance or vehicle controllers.

\begin{definitionbox}{Microprocessor}
A central processing unit implemented on a single chip; a general-purpose
device that needs external memory and I/O.
\end{definitionbox}

\begin{definitionbox}{Microcontroller}
A microprocessor together with integrated memory and input/output peripherals
on one chip, used in embedded and special-purpose systems.
\end{definitionbox}

\begin{center}
\begin{tabular}{p{0.28\textwidth}p{0.30\textwidth}p{0.30\textwidth}}
\toprule
\textbf{Feature} & \textbf{Microprocessor} & \textbf{Microcontroller}\\
\midrule
Contents & CPU on a chip & CPU + memory + I/O on a chip\\
Purpose & General-purpose & Embedded / special-purpose\\
System & Needs external parts & Self-contained on one chip\\
Example & EX-01 office desktop & Appliance controller\\
\bottomrule
\end{tabular}
\end{center}

The key difference is the \textbf{level of integration}, not clock speed. A
microcontroller is not ``just a microprocessor'': it is a more integrated
device that already includes memory and I/O on the same chip.

\section{Diagnosing a ``slow'' computer (EX-01)}
The office desktop in this running example is a quad-core 2.4\,GHz machine
with 16\,GB of RAM and an SSD. Suppose one particular program still runs
slowly. A common but incorrect conclusion is that the CPU must be slow,
because ``more cores should mean faster''. More cores do not help a
single-threaded program, which uses only one core at a time. The real cause is
usually a \textbf{bottleneck}: memory pressure, disk input/output, or software
that is not multithreaded.

\begin{example}
A question describes a quad-core 2.4\,GHz desktop with 16\,GB RAM and an SSD,
and asks why one application is slow. The best answer is to diagnose a stated
bottleneck --- a single-threaded program, insufficient memory, or slow disk
input/output --- rather than to declare the CPU slow. A slow program is not
automatically a ``slow CPU.''
\end{example}

\section{Exam retrieval and common confusions}
Use two retrieval directions. Given the processor, recall its core profile;
given the workload, recall the better processor.
\begin{itemize}
  \item CPU $\leftrightarrow$ few powerful general-purpose cores, low latency.
  \item GPU $\leftrightarrow$ many simple cores, high throughput, parallel.
  \item Serial $\leftrightarrow$ one task at a time in sequence.
  \item Parallel $\leftrightarrow$ several tasks or operations at once.
  \item Microprocessor $\leftrightarrow$ CPU on a single chip, general-purpose.
  \item Microcontroller $\leftrightarrow$ microprocessor + memory + I/O on one
    chip, embedded.
  \item Cache order $\leftrightarrow$ L1 (smallest, fastest) to L3 (largest,
    slowest).
\end{itemize}

Three distinctions prevent most errors in this topic:
\begin{enumerate}
  \item Do not treat GPU cores as general-purpose CPU cores. They are simpler
    and optimised for parallel throughput.
  \item Do not confuse serial with parallel. Serial is sequential, one task at
    a time; parallel is simultaneous.
  \item Do not attribute a slow program on a capable machine automatically to
    a slow CPU. Check the multithreading, memory, and disk-I/O bottleneck; and
    do not confuse a microprocessor (CPU on a chip) with a microcontroller
    (CPU + memory + I/O on one chip).
\end{enumerate}

\subsection*{Exercises}
\begin{enumerate}[label=\arabic*.]
  \item Distinguish serial processing from parallel processing.
  \item State one design goal of the CPU and one of the GPU, and give the
    core profile of each.
  \item Explain the difference between a microprocessor and a microcontroller.
  \item What does hyper-threading allow, and why is it not the same as adding
    a physical core?
  \item A quad-core 2.4\,GHz desktop with 16\,GB RAM and an SSD runs one
    program slowly. Give the most appropriate diagnosis and say why.
\end{enumerate}

\subsection*{Solutions}
\begingroup\small
\begin{enumerate}[label=\arabic*.]
  \item Serial processing performs one task at a time in sequence; parallel
    processing performs several tasks or operations at the same time.
  \item The CPU is optimised for low latency and has few powerful
    general-purpose cores. The GPU is optimised for high throughput and has
    many simpler cores for parallel work.
  \item A microprocessor is a CPU on a single chip and is general-purpose,
    needing external memory and I/O. A microcontroller adds memory and I/O
    peripherals on the same chip, forming an embedded, special-purpose system;
    the difference is the level of integration.
  \item Hyper-threading lets one physical core present two logical threads, so
    it can stay busy while one thread waits. It does not add a physical core;
    a quad-core chip with hyper-threading still has four physical cores.
  \item The program is likely limited by a bottleneck --- a single-threaded
    workload, insufficient memory, or slow disk input/output --- rather than
    by a slow CPU, since extra cores do not speed up a single-threaded
    program. A slow program is not automatically a ``slow CPU.''
\end{enumerate}
\endgroup
\end{document}
